\documentclass[10pt,a4paper]{amsart}
\usepackage{amsmath,amssymb,amsthm,graphicx}
% For commutative diagrams we use Xy-pic
\usepackage[all]{xy}

% We use 2cell for 2-commutative diagrams.
\xyoption{2cell}
% Theorem environments.
%
\theoremstyle{plain}
\newtheorem{theorem}[subsection]{Theorem}
\newtheorem{proposition}[subsection]{Proposition}
\newtheorem{lemma}[subsection]{Lemma}

\theoremstyle{definition}
\newtheorem{definition}[subsection]{Definition}
\newtheorem{example}[subsection]{Example}
\newtheorem{exercise}[subsection]{Exercise}
\newtheorem{situation}[subsection]{Situation}

\theoremstyle{remark}
\newtheorem{remark}[subsection]{Remark}
\newtheorem{remarks}[subsection]{Remarks}

\numberwithin{equation}{subsection}

% Macros
%
\def\lim{\mathop{\mathrm{lim}}\nolimits}
\def\colim{\mathop{\mathrm{colim}}\nolimits}
\def\Spec{\mathop{\mathrm{Spec}}}
\def\Hom{\mathop{\mathrm{Hom}}\nolimits}
\def\Ext{\mathop{\mathrm{Ext}}\nolimits}
\def\SheafHom{\mathop{\mathcal{H}\!\mathit{om}}\nolimits}
\def\SheafExt{\mathop{\mathcal{E}\!\mathit{xt}}\nolimits}
\def\Sch{\mathit{Sch}}
\def\Mor{\mathop{\mathrm{Mor}}\nolimits}
\def\Ob{\mathop{\mathrm{Ob}}\nolimits}
\def\Sh{\mathop{\mathit{Sh}}\nolimits}
\def\NL{\mathop{N\!L}\nolimits}
\def\CH{\mathop{\mathrm{CH}}\nolimits}
\def\proetale{{pro\text{-}\acute{e}tale}}
\def\etale{{\acute{e}tale}}
\def\QCoh{\mathit{QCoh}}
\def\Ker{\mathop{\mathrm{Ker}}}
\def\Im{\mathop{\mathrm{Im}}}
\def\Coker{\mathop{\mathrm{Coker}}}
\def\Coim{\mathop{\mathrm{Coim}}}

% Boxtimes
%
\DeclareMathSymbol{\boxtimes}{\mathbin}{AMSa}{"02}

%
% Macros for moduli stacks/spaces
%
\def\QCohstack{\mathcal{QC}\!\mathit{oh}}
\def\Cohstack{\mathcal{C}\!\mathit{oh}}
\def\Spacesstack{\mathcal{S}\!\mathit{paces}}
\def\Quotfunctor{\mathrm{Quot}}
\def\Hilbfunctor{\mathrm{Hilb}}
\def\Curvesstack{\mathcal{C}\!\mathit{urves}}
\def\Polarizedstack{\mathcal{P}\!\mathit{olarized}}
\def\Complexesstack{\mathcal{C}\!\mathit{omplexes}}
% \Pic is the operator that assigns to X its picard group, usage \Pic(X)
% \Picardstack_{X/B} denotes the Picard stack of X over B
% \Picardfunctor_{X/B} denotes the Picard functor of X over B
\def\Pic{\mathop{\mathrm{Pic}}\nolimits}
\def\Picardstack{\mathcal{P}\!\mathit{ic}}
\def\Picardfunctor{\mathrm{Pic}}
\def\Deformationcategory{\mathcal{D}\!\mathit{ef}}
\usepackage{cleveref}
\crefname{originaldefinition}{original definition}{original definitions}
\Crefname{originaldefinition}{Original definition}{Original definitions}
\makeatletter
\newlabel{definition-cech-complex}{{03OL}{source}{}{sourcepointer-1}{}}
\newlabel{definition-cech-complex@cref}{{[originaldefinition][0][]03OL}{[1][0][]source}{}{}{}}
\makeatother
\title{Cech cohomology of an injective presheaf vanishes in positive degrees}
\author{The Stacks Project Authors}
\date{stacks-project:etale-cohomology.tex:local-snapshot-20261011; GNU FDL 1.2 or later}
\begin{document}
\maketitle
\noindent\textit{Editorial note.} All delivered mathematical text is the author's own, copied from the pinned Stacks Project chapter file etale-cohomology.tex; no statement or proof step was rewritten, reordered or removed. The check accent on Cech cohomology is reproduced as the author's own \textbackslash{}check H notation.\par\medskip
\noindent\textit{Editorial note.} The preamble is the chapter's own native preamble excerpt used for this source: the Xy-pic setup at lines 14-18 and the theorem declarations with the mathematical macros at lines 154-217 of sources/stacks-etale-cohomology/preamble.tex. The quoted class line and the xr-hyper externaldocument block are not included because the assembler supplies a plain amsart wrapper; no macro was invented or redefined.\par\medskip
\noindent\textit{Allowed standard prerequisites (editorial).}\begin{itemize}
\item Cech complex and Cech cohomology of a family in a presheaf of abelian groups. Same chapter, Definition environment labelled definition-cech-complex (source line 1477, Stacks tag 03OL); the statement picks up its notation and assumptions. Supplies the complex Z\textasciicircum{}bullet\_U and the groups H\textasciicircum{}p-check; the vanishing argument is the counted content.
\item The complex Z\textasciicircum{}bullet\_U is exact in negative degrees and Hom(-, I) into an injective is exact. Same chapter, the immediately preceding computation of the complex Z\textasciicircum{}bullet\_U (the author's phrase 'we have just seen'), together with the defining property of an injective object. Supplies the two facts combined in the proof; the combination producing exactness in positive degrees is retained in the counted span.
\end{itemize}\medskip
\ifdefined\hypertarget\hypertarget{sourcepointer-1}{}\fi
\noindent\textit{Source reference.} Original definition 03OL: sources/stacks-etale-cohomology/pinned-original-source.tex, source lines 1477--1477. Same-chapter pointer named in the first line of the counted statement (line 1703) as the source of notation and assumptions; the identification of the Cech complex as a Hom complex and the exactness argument are retained verbatim in the counted span.\par\medskip
\begin{lemma}
\label{lemma-hom-injective}
Notation and assumptions as in Definition \ref{definition-cech-complex}.
If $\mathcal{I}$ is an injective object of $\textit{PAb}(\mathcal{C})$,
then $\check H^p(\mathcal{U}, \mathcal{I}) = 0$ for all $p > 0$.
\end{lemma}

\begin{proof}
The {\v C}ech complex is the result of applying the functor
$\Hom_{\textit{PAb}(\mathcal{C})}(-, \mathcal{I}) $ to the complex $
\mathbf{Z}^\bullet_\mathcal{U} $, i.e.,
$$
\check H^p(\mathcal{U}, \mathcal{I}) = H^p
(\Hom_{\textit{PAb}(\mathcal{C})} (\mathbf{Z}^\bullet_\mathcal{U},
\mathcal{I})).
$$
But we have just seen that $\mathbf{Z}^\bullet_\mathcal{U}$ is exact in
negative degrees, and the functor $\Hom_{\textit{PAb}(\mathcal{C})}(-,
\mathcal{I})$ is exact, hence $\Hom_{\textit{PAb}(\mathcal{C})}
(\mathbf{Z}^\bullet_\mathcal{U}, \mathcal{I})$ is exact in positive degrees.
\end{proof}

\end{document}
